\documentclass[10pt,a4paper,twocolumn]{article}%
\usepackage[T1]{fontenc}%
\usepackage[utf8]{inputenc}%
\usepackage{lmodern}%
\usepackage{textcomp}%
\usepackage{lastpage}%
\usepackage{graphicx}%
\usepackage{amsmath}%
\usepackage{amssymb}%
\usepackage{booktabs}%
\usepackage{hyperref}%
%
%
%
\begin{document}%
\normalsize%
\author{Birrul Wildan}%
\title{A More Complex PyLaTeX Document Example}%
\date{\today}%
\maketitle%
\section{Abstract}%
\label{sec:Abstract}%
This document demonstrates a more complex usage of PyLaTeX, including sections, subsections, figures, tables, lists, and mathematical equations. It serves as an example for generating structured LaTeX documents programmatically.

%
\section{Introduction}%
\label{sec:Introduction}%
Welcome to this advanced PyLaTeX example. We will explore various features to create a rich document.%
\subsection{Document Structure}%
\label{subsec:DocumentStructure}%
LaTeX documents are structured using sections and subsections.

%
\section{Figures}%
\label{sec:Figures}%
Here is an example of including a figure. Please ensure you have an image file named `example{-}image{-}a.png` in the same directory for this to compile correctly, or replace it with an existing image path.%


\begin{figure}[h!]%
\centering%
\includegraphics[width=0.8\textwidth]{example-image-a.png}%
\caption{A placeholder image demonstrating figure inclusion.}%
\label{fig:example\_image}%
\end{figure}

%
\section{Tables}%
\label{sec:Tables}%
Below is an example of a simple table.%
\begin{tabular}{c|c|c}%
\hline%
Header 1&Header 2&Header 3\\%
\hline%
Row 1 Col 1&Row 1 Col 2&Row 1 Col 3\\%
Row 2 Col 1&Row 2 Col 2&Row 2 Col 3\\%
\hline%
\end{tabular}%
\label{tab:example\_table}

%
\section{Lists}%
\label{sec:Lists}%
Here are some examples of lists:%
\begin{itemize}%
\item%
First item in the list.%
\item%
Second item, which is a bit longer to show wrapping.%
\item%
Third item.%
\end{itemize}

%
\section{Mathematics}%
\label{sec:Mathematics}%
We can include inline math like \$E=mc\^{}2\$ or displayed equations:%
\[ \alpha + \beta = \gamma \]%
And aligned equations:%
\begin{alignat}{2}%
f(x) &= x^2 + 2x + 1 \\%
     &= (x+1)^2%
\end{alignat}%
\label{eq:quadratic}

%
\section{Conclusion}%
\label{sec:Conclusion}%
This document has showcased various features of PyLaTeX. You can refer to Figure 
ef\{fig:example\_image\} and Table 
ef\{tab:example\_table\}. Equation 
ef\{eq:quadratic\} is an example of aligned math.%
\newline%
\textit{Happy LaTeXing with Python!}

%
\end{document}